% Advanced OmniLaTeX Beamer Showcase
% Compile with: latexmk -lualatex main.tex
\documentclass[doctype=beamer, beameraspectratio=169, beamertheme=Madrid, institution=generic]{omnilatex}
\usecolortheme{whale}
\usepackage{tikz}
\usetikzlibrary{arrows.meta,positioning,calc,fit}
\usepackage{listings,xcolor,amsmath,amssymb,booktabs}
\usepackage{tcolorbox}
\tcbuselibrary{skins,breakable}
\lstset{
  basicstyle=\ttfamily\small,
  keywordstyle=\color{blue}\bfseries,
  commentstyle=\color{green!50!black}\itshape,
  stringstyle=\color{red!70!black},
  numbers=left, numberstyle=\tiny\color{gray},
  frame=single, breaklines=true, showstringspaces=false, tabsize=2,
}
\newcommand{\checkmark}{\text{\(\checkmark\)}}
\title[Advanced Beamer]{Advanced Beamer Features\\with OmniLaTeX and LuaLaTeX}
\subtitle{A Comprehensive Showcase}
\author[Jane Doe]{Jane Doe}
\institute[OmniLaTeX Project]{OmniLaTeX Project\\\texttt{github.com/WyattAu/OmniLaTeX-template}}
\date{June 2026}
\usepackage[style=ieee]{biblatex}
\addbibresource{../../bib/bibliography.bib}
\begin{document}
\begin{frame}[plain]\titlepage\end{frame}
\begin{frame}{Outline}
  \tableofcontents[hideallsubsections]
\end{frame}
\section{Overlays and Animations}
\begin{frame}{Conditional Content}
  \begin{block}{\texttt{\textbackslash only} vs.\ \texttt{\textbackslash uncover} vs.\ \texttt{\textbackslash visible}}
    \begin{itemize}
      \item \texttt{\textbackslash only<2->{}}: Appears from slide 2; space reserved only when visible.
      \item \texttt{\textbackslash uncover<2-3>{}}: Visible on slides 2--3; space always reserved.
      \item \texttt{\textbackslash visible<2->{}}: Like \texttt{\textbackslash uncover}, hidden content not in PDF.
    \end{itemize}
  \end{block}
  \only<2->{\begin{alertblock}{This block appears on slide 2 onward}
    Demonstrates \texttt{\textbackslash only<2->{}}.
  \end{alertblock}}
  \uncover<3-4>{\begin{exampleblock}{This block is visible on slides 3--4}
    Demonstrates \texttt{\textbackslash uncover<3-4>{}}.
  \end{exampleblock}}
\end{frame}

\begin{frame}{Step-by-Step Reveals with \texttt{\textbackslash pause}}
  Step 1: First piece of content.\pause
  Step 2: Second piece of content.\pause
  Step 3: Third piece of content.\pause
  Step 4: Final piece of content.
  \vspace{1em}
  \begin{block}{Note}
    \texttt{\textbackslash pause} inserts a slide break at each invocation.
  \end{block}
\end{frame}

\begin{frame}{Switching, Temporal, and Incremental Items}
  \begin{columns}[T]
    \column{0.48\textwidth}
    \begin{block}{\texttt{\textbackslash alt}}
      \alt<1|2>{\textbf{Option A: Procedural}}{\textbf{Option B: Functional}}
    \end{block}
    \begin{block}{\texttt{\textbackslash temporal}}
      \temporal<1|2|3>{Before}{Processing \textbf{now}}{After}
    \end{block}
    \column{0.48\textwidth}
    \begin{block}{Incremental Items}
      \begin{itemize}
        \item<1-> First item: always visible.
        \item<2-> Second item: appears on slide 2.
        \item<3-> Third item: appears on slide 3.
      \end{itemize}
    \end{block}
  \end{columns}
  \vspace{1em}
  \begin{block}{Custom Bullets and Highlighted Emphasis}
    \begin{columns}[T]
      \column{0.48\textwidth}
      \begin{itemize}
        \item[\checkmark] Completed task
        \item[$\times$] Failed task
        \item[$\diamond$] Optional task
      \end{itemize}
      \column{0.48\textwidth}
      \begin{center}
        \only<2>{\color{red}\textbf{Highlighted on slide 2 only}}
      \end{center}
    \end{columns}
  \end{block}
\end{frame}
\section{Blocks and Environments}
\begin{frame}{Block Environments}
  \begin{columns}[T]
    \column{0.48\textwidth}
    \begin{block}{Standard Block}General-purpose highlighted content.\end{block}
    \begin{alertblock}{Alert Block}Warnings and critical information.\end{alertblock}
    \column{0.48\textwidth}
    \begin{exampleblock}{Example Block}Demonstrations and illustrative content.\end{exampleblock}
    \begin{block}{Custom Colored Block}
      \setbeamercolor{block body}{bg=blue!10}
      Adapt block colors per-block.
    \end{block}
  \end{columns}
\end{frame}

\begin{frame}{Blocks with Overlays}
  \begin{block}<2->{Definition Block}
    This block appears only from slide 2 onward.
  \end{block}
  \only<3>{\begin{alertblock}{Alert on Slide 3 Only}
    This warning appears only on the third sub-slide.
  \end{alertblock}}
  \uncover<4->{\begin{exampleblock}{Example Block (Slide 4+)}
    This example block becomes visible from slide 4.
  \end{exampleblock}}
\end{frame}
\section{Columns and Layouts}
\begin{frame}{Side-by-Side Comparison}
  \begin{columns}[T]
    \column{0.48\textwidth}
    \begin{block}{Approach A: Procedural}
      \begin{enumerate}
        \item Read input file
        \item Parse line by line
        \item Transform tokens
        \item Write output file
      \end{enumerate}
    \end{block}
    \column{0.48\textwidth}
    \begin{block}{Approach B: Declarative}
      \begin{enumerate}
        \item Define grammar rules
        \item Specify transformations
        \item Apply pipeline
        \item Collect results
      \end{enumerate}
    \end{block}
  \end{columns}
  \vspace{0.5em}
  \rule{1\textwidth}{0.4pt}
  \vspace{0.5em}
  \begin{minipage}{0.3\textwidth}
    \begin{block}{Phase 1: Input}Parse and validate parameters.\end{block}
  \end{minipage}\hfill
  \begin{minipage}{0.3\textwidth}
    \begin{block}{Phase 2: Process}Execute the core algorithm.\end{block}
  \end{minipage}\hfill
  \begin{minipage}{0.3\textwidth}
    \begin{block}{Phase 3: Output}Format and produce results.\end{block}
  \end{minipage}
\end{frame}
\section{TikZ Integration}
\begin{frame}{Animated Flowchart Construction}
  \begin{center}
    \begin{tikzpicture}[
      node distance=1.5cm,
      block/.style={rectangle, draw, fill=blue!5, minimum width=2cm, minimum height=0.8cm, align=center, rounded corners},
      highlight/.style={rectangle, draw, fill=red!20, minimum width=2cm, minimum height=0.8cm, align=center, rounded corners, thick},
    ]
      \node[block] (input) {Input};
      \only<2->{\node[block, right=of input] (process) {Process};}
      \only<3->{\node[block, right=of process] (decision) {Decision};}
      \only<4->{\node[block, right=of decision] (output) {Output};}
      \only<2->{\draw[->,>=stealth,thick] (input) -- (process);}
      \only<3->{\draw[->,>=stealth,thick] (process) -- (decision);}
      \only<4->{\draw[->,>=stealth,thick] (decision) -- (output);}
      \only<3>{\node[highlight, right=of process] (decision) {Decision};}
    \end{tikzpicture}
  \end{center}
  \begin{center}
    \small Step-by-step flowchart construction using \texttt{\textbackslash only}
  \end{center}
\end{frame}

\begin{frame}{Step-by-Step Algorithm Visualization}
  \begin{center}
    \begin{tikzpicture}[
      node distance=1.2cm,
      state/.style={circle, draw, minimum size=0.9cm, font=\small},
      visited/.style={state, fill=green!20, draw=green!60!black},
      current/.style={state, fill=yellow!40, draw=orange, thick},
    ]
      \node[state] (s) at (0,0) {$s$};
      \node[state] (a) at (2,1.5) {$a$};
      \node[state] (b) at (2,-1.5) {$b$};
      \node[state] (c) at (4,0) {$c$};
      \node[state] (d) at (6,1.5) {$d$};
      \node[state] (t) at (6,-1.5) {$t$};
      \draw[->,>=stealth] (s) -- (a) node[midway,above left] {3};
      \draw[->,>=stealth] (s) -- (b) node[midway,below left] {5};
      \draw[->,>=stealth] (a) -- (c) node[midway,above] {2};
      \draw[->,>=stealth] (b) -- (c) node[midway,below] {1};
      \draw[->,>=stealth] (c) -- (d) node[midway,above right] {4};
      \draw[->,>=stealth] (c) -- (t) node[midway,below right] {6};
      \draw[->,>=stealth] (d) -- (t) node[midway,right] {2};
      \only<2>{\node[visited] at (0,0) {$s$};}
      \only<3>{\node[visited] at (0,0) {$s$}; \node[visited] at (2,1.5) {$a$};}
      \only<4>{\node[visited] at (0,0) {$s$}; \node[visited] at (2,1.5) {$a$}; \node[visited] at (2,-1.5) {$b$};}
      \only<5>{\node[visited] at (0,0) {$s$}; \node[visited] at (2,1.5) {$a$}; \node[visited] at (2,-1.5) {$b$}; \node[current] at (4,0) {$c$};}
      \only<6>{\node[visited] at (0,0) {$s$}; \node[visited] at (2,1.5) {$a$}; \node[visited] at (2,-1.5) {$b$}; \node[visited] at (4,0) {$c$}; \node[visited] at (6,1.5) {$d$};}
      \only<7>{\node[visited] at (0,0) {$s$}; \node[visited] at (2,1.5) {$a$}; \node[visited] at (2,-1.5) {$b$}; \node[visited] at (4,0) {$c$}; \node[visited] at (6,1.5) {$d$}; \node[visited] at (6,-1.5) {$t$};}
    \end{tikzpicture}
  \end{center}
  \small Dijkstra's algorithm: visited (green), current (yellow)
\end{frame}

\begin{frame}{Highlighted Subgraphs}
  \begin{center}
    \begin{tikzpicture}[
      every node/.style={circle, draw, minimum size=1cm},
      subgraph/.style={draw, dashed, thick, rounded corners, inner sep=5pt},
    ]
      \node (a) at (0,0) {$A$};
      \node (b) at (2,0) {$B$};
      \node (c) at (1,1.73) {$C$};
      \node (d) at (4,0) {$D$};
      \node (e) at (5,1.73) {$E$};
      \node (f) at (6,0) {$F$};
      \draw[->,>=stealth] (a) -- (b);
      \draw[->,>=stealth] (b) -- (c);
      \draw[->,>=stealth] (c) -- (a);
      \draw[->,>=stealth] (b) -- (d);
      \draw[->,>=stealth] (d) -- (e);
      \draw[->,>=stealth] (e) -- (f);
      \draw[->,>=stealth] (f) -- (d);
      \only<2>{\node[subgraph, fill=blue!10, fit=(a)(b)(c)] {};
                \node[below=0.3cm of a, blue!70!black, font=\small\bfseries] {Subgraph 1};}
      \only<3>{\node[subgraph, fill=orange!10, fit=(d)(e)(f)] {};
                \node[below=0.3cm of f, orange!70!black, font=\small\bfseries] {Subgraph 2};}
      \only<4>{\node[subgraph, fill=blue!10, fit=(a)(b)(c)] {};
                \node[subgraph, fill=orange!10, fit=(d)(e)(f)] {};
                \node[below=0.3cm of a, blue!70!black, font=\small\bfseries] {Subgraph 1};
                \node[below=0.3cm of f, orange!70!black, font=\small\bfseries] {Subgraph 2};}
    \end{tikzpicture}
  \end{center}
  \begin{center}
    \small Highlight subgraphs with \texttt{\textbackslash only} and TikZ \texttt{fit}
  \end{center}
\end{frame}
\section{Math in Beamer}
\begin{frame}{Step-by-Step Equation Derivation}
  \onslide<2->{%
    \begin{equation}
      \int_{-\infty}^{\infty} e^{-x^2} \, dx = \sqrt{\pi}
    \end{equation}
  }
  \vspace{1em}
  \onslide<3->{%
    \begin{block}{Gaussian Integral --- Derivation}
      \[
        I = \int_{-\infty}^{\infty} e^{-x^2} \, dx
        \implies I^2 = \pi
        \implies I = \sqrt{\pi}
      \]
    \end{block}
  }
\end{frame}

\begin{frame}{Matrix with Highlighted Elements}
  \[
    A = \begin{pmatrix}
      \only<2>{\color{red}\mathbf{a_{11}}} \only<3->{a_{11}} &
      \only<3>{\color{blue}\mathbf{a_{12}}} \only<4->{a_{12}} &
      a_{13} \\
      a_{21} &
      \only<4>{\color{green!50!black}\mathbf{a_{22}}} \only<5->{a_{22}} &
      a_{23} \\
      a_{31} & a_{32} &
      \only<5>{\color{purple}\mathbf{a_{33}}} \only<6->{a_{33}}
    \end{pmatrix}
  \]
  \vspace{1em}
  \onslide<6->
  \begin{block}{Determinant Expansion}
    \[
      \det(A) = a_{11}(a_{22}a_{33} - a_{23}a_{32})
              - a_{12}(a_{21}a_{33} - a_{23}a_{31})
              + a_{13}(a_{21}a_{32} - a_{22}a_{31})
    \]
  \end{block}
\end{frame}

\begin{frame}{Proof Steps with Overlays}
  \begin{block}{Theorem: Pythagorean Identity}
    For any angle $\theta$: $\sin^2\theta + \cos^2\theta = 1$
  \end{block}
  \vspace{1em}
  \begin{proof}
    \only<2>{Start with the unit circle definition.}
    \only<3>{For any point $(x, y)$ on the unit circle: $x^2 + y^2 = 1$.}
    \only<4>{Substitute $x = \cos\theta$ and $y = \sin\theta$:}
    \only<5>{We obtain $\cos^2\theta + \sin^2\theta = 1$. \qedhere}
  \end{proof}
\end{frame}
\section{Code in Beamer}
\begin{frame}[fragile]{Python with Line Highlighting}
  \begin{block}{Normalized Function}
    \begin{lstlisting}[language=Python, highlightlines={3-5}]
import numpy as np

def normalize(x):
    """Normalize array to [0, 1] range."""
    x_min, x_max = x.min(), x.max()
    return (x - x_min) / (x_max - x_min)

data = np.random.rand(100)
result = normalize(data)
    \end{lstlisting}
  \end{block}
\end{frame}
\begin{frame}[fragile]{Rust Code Listing}
  \begin{block}{Rust: Iterator Adapter}
    \begin{lstlisting}[language=Rust]
fn fibonacci() -> impl Iterator<Item = u64> {
    let mut a: u64 = 0;
    let mut b: u64 = 1;
    std::iter::from_fn(move || {
        let next = a;
        let sum = a + b;
        a = b;
        b = sum;
        Some(next)
    })
}

fn main() {
    let fibs: Vec<u64> = fibonacci().take(10).collect();
    println!("{:?}", fibs);
}
    \end{lstlisting}
  \end{block}
\end{frame}
\section{Tables in Beamer}
\begin{frame}{Comparison Tables with Highlighted Rows}
  \begin{center}
    \begin{tabular}{l c c c}
      \toprule
      \textbf{Feature} & \textbf{Option A} & \textbf{Option B} & \textbf{Option C} \\
      \midrule
      Speed       & High  & Medium & Low  \\
      \rowcolor{green!10} Accuracy  & High  & High   & Medium \\
      Memory      & Low   & Medium & High \\
      \rowcolor{red!10} Complexity & Low   & High   & High   \\
      Scalability & Medium & High  & High \\
      \bottomrule
    \end{tabular}
  \end{center}
  \vspace{1em}
  \begin{exampleblock}{Highlighting Rows}
    Colored rows with \texttt{\textbackslash rowcolor} draw attention to
    key comparisons: green for strengths, red for weaknesses.
  \end{exampleblock}
\end{frame}
\section{Figures}
\begin{frame}{Subfigures with Overlays}
  \begin{columns}[T]
    \column{0.48\textwidth}
    \begin{block}{Figure 1a: Linear Growth}
      \begin{center}
        \begin{tikzpicture}[scale=0.7]
          \begin{axis}[width=5.5cm, height=4cm, xlabel={Steps}, ylabel={Value}, grid=major, thick]
            \addplot[blue, thick, smooth, domain=0:10, samples=20] {0.5*x};
          \end{axis}
        \end{tikzpicture}
      \end{center}
      \centering\small Linear: $f(x) = 0.5x$
    \end{block}
    \column{0.48\textwidth}
    \begin{block}<2->{Figure 1b: Exponential Growth}
      \begin{center}
        \begin{tikzpicture}[scale=0.7]
          \begin{axis}[width=5.5cm, height=4cm, xlabel={Steps}, ylabel={Value}, grid=major, thick]
            \addplot[red, thick, smooth, domain=0:10, samples=30] {exp(0.3*x)};
          \end{axis}
        \end{tikzpicture}
      \end{center}
      \centering\small Exponential: $f(x) = e^{0.3x}$
    \end{block}
  \end{columns}
  \begin{center}
    \small Second subfigure appears on overlay 2, demonstrating timed figure reveals.
  \end{center}
\end{frame}
\section{Navigation and Structure}
\begin{frame}{Table of Contents}
  \begin{columns}[T]
    \column{0.5\textwidth}
    \begin{block}{Full Table of Contents}
      \tableofcontents[hideallsubsections]
    \end{block}
    \column{0.5\textwidth}
    \begin{block}{Current Section Only}
      \tableofcontents[currentsection, hideallsubsections]
    \end{block}
  \end{columns}
\end{frame}

\subsection{Subsection Demonstration}
\begin{frame}{Subsection Content}
  This frame belongs to a subsection, demonstrating hierarchical structure.
  \begin{block}{Key Point}
    Use \texttt{\textbackslash section} and \texttt{\textbackslash subsection}
    to organize your presentation.
  \end{block}
\end{frame}
\section{Framed Content}
\begin{frame}{tcolorbox with Beamer Overlays}
  \begin{tcolorbox}[
    enhanced, colback=blue!5, colframe=blue!50!black,
    title={Overlay-Aware tcolorbox}, fonttitle=\bfseries,
    borderline west={2.5pt}{0pt}{blue!50!black},
  ]
    This tcolorbox is always visible. It supports full beamer overlay
    syntax because tcolorbox integrates deeply with beamer.
  \end{tcolorbox}
  \vspace{1em}
  \only<2>{%
    \begin{tcolorbox}[
      enhanced, colback=red!5, colframe=red!70!black,
      title={Slide 2 Only}, fonttitle=\bfseries,
      borderline west={2.5pt}{0pt}{red!70!black},
    ]
      This tcolorbox appears only on slide 2.
    \end{tcolorbox}
  }
  \only<3>{%
    \begin{tcolorbox}[
      enhanced, colback=green!5, colframe=green!50!black,
      title={Slide 3 Only}, fonttitle=\bfseries,
      borderline west={2.5pt}{0pt}{green!50!black},
    ]
      This tcolorbox appears only on slide 3.
    \end{tcolorbox}
  }
\end{frame}
\section{Footnotes and References}
\begin{frame}{Footnotes and Citations}
  \begin{block}{Timed Footnotes}
    This content is always visible.\footnote<2->{This footnote appears on slide 2 onward.}
    Additional context.\footnote<3->{This footnote appears on slide 3 onward.}
  \end{block}
  \vspace{1em}
  \begin{block}{Citations}
    \begin{itemize}
      \item Einstein's special relativity\cite{einstein1905}.
      \item Dirac's quantum mechanics\cite{diracPrinciplesQuantumMechanics1981}.
      \item Bird et al.\ on NLP with Python\cite{birdNaturalLanguageProcessing2009}.
    \end{itemize}
  \end{block}
  \vspace{1em}
  \begin{center}
    \small
    \printbibliography[heading=none]
  \end{center}
\end{frame}
%% ════════════════════════════════════════════════════
%% SUMMARY & THANK YOU
%% ════════════════════════════════════════════════════

\begin{frame}{Summary}
  \begin{enumerate}
    \item Overlays enable step-by-step content delivery.
    \item Blocks provide structured, colored content areas.
    \item Columns support flexible multi-column layouts.
    \item TikZ diagrams animate and highlight per-slide.
    \item Math supports progressive equation building.
    \item Code listings with line highlighting.
    \item Tables with row highlighting for comparisons.
    \item Subfigures with timed reveals.
    \item tcolorbox integrates with beamer overlays.
    \item Footnotes, citations, and appendix complete the toolkit.
  \end{enumerate}
\end{frame}

\begin{frame}[plain,c]
  \vspace{2cm}
  {\huge\bfseries Thank You!}\\[1cm]
  {\large Questions?}\\[2cm]
  {\footnotesize
  \texttt{https://github.com/WyattAu/OmniLaTeX-template}\\
  \texttt{omnilatex-template.wyattau.com}
  }
\end{frame}

%% ════════════════════════════════════════════════════
%% APPENDIX
%% ════════════════════════════════════════════════════

\appendix

\begin{frame}[allowframebreaks]{Appendix: Full Doctype Reference}
  \begin{table}
    \centering
    \begin{tabular}{l l l}
      \hline
      Doctype         & Base Class & Aliases \\
      \hline
      thesis          & scrbook    & theses \\
      dissertation    & scrbook    & dissertations \\
      book            & scrbook    & --- \\
      manual          & scrreprt   & manuals, guide, handbook \\
      article         & scrartcl   & articles, paper \\
      cv              & scrartcl   & resume, curriculumvitae \\
      poster          & scrartcl   & posters \\
      beamer          & beamer     & --- \\
      \hline
    \end{tabular}
  \end{table}
  (27 canonical profiles with 55+ aliases)
\end{frame}

\begin{frame}{Appendix: Detailed Derivation}
  \begin{block}{Cauchy--Schwarz Inequality Proof Sketch}
    For $u, v \in V$ over field $\mathbb{F}$:
    \[ t \mapsto \|u + tv\|^2 = \langle u + tv, u + tv \rangle \geq 0 \]
    Expanding:
    \[ \|u\|^2 + 2t \operatorname{Re}\langle u, v \rangle + t^2 \|v\|^2 \geq 0 \]
    This quadratic in $t$ is non-negative, so its discriminant satisfies:
    \[ \Delta = 4(\operatorname{Re}\langle u, v \rangle)^2 - 4\|u\|^2\|v\|^2 \leq 0 \]
    Therefore $|\langle u, v \rangle|^2 \leq \|u\|^2 \|v\|^2$.
  \end{block}
\end{frame}

\end{document}
